\clearpage\thispagestyle{empty}
\pdfbookmark[0]{Sources for Calculation of Easter (1800)}{osterfest1800-sources}
\begin{center}\textbf{Sources for Calculation of Easter (1800)}\end{center}
\noindent This AI translation follows Gauss’s German text in \textit{Werke}, vol.~VI (Göttingen, 1874), pp.~73--79, including the general rules and the handwritten remark printed at the end. The first printing appeared in Franz Xaver von Zach’s \textit{Monatliche Correspondenz zur Beförderung der Erd- und Himmels-Kunde}, vol.~2, August 1800, article~XV, pp.~121--130, under the name Doctor Gaus in Braunschweig. \href{https://www.digitale-sammlungen.de/de/view/bsb10538594?page=127}{First printing}. \href{https://gdz.sub.uni-goettingen.de/id/PPN236007467}{Collected German text}.
\par\smallskip\noindent The first printing gives 9 possible values of $a$ on p.~123 and lists 0 through 18 across pp.~123--124. The collected note gives 19. The first printing expressly excludes 1700 when counting leap years (p.~126). The collected text specifies inclusion of the given year if it is a leap year. For the Gregorian calendar through 1699, the first printing gives $M=12$, $N=2$ (p.~128). The collected text gives $M=22$, $N=2$ (p.~78).
\par\smallskip\noindent A short English rendering of the three initial instructions and result appears in Eric S. Roberts’s \textit{Programming Abstractions in C++}, Figure~2-18, \href{https://tfetimes.com/wp-content/uploads/2015/04/Prog-in-Cpp.pdf}{Spring 2012}, p.~119, and \href{https://web.stanford.edu/dept/cs_edu/resources/textbook/Reader-Beta-2012.pdf}{Autumn 2012}, pp.~118--119 (exercise~8 and figure). For calendar terminology and the distinct later rules, see Reinhold Bien, \href{https://doi.org/10.1007/s00407-004-0078-5}{\textit{Gauß and Beyond: The Making of Easter Algorithms}}, \textit{Archive for History of Exact Sciences} 58 (2004), pp.~439--452, especially pp.~439--442.
\par\smallskip\noindent Translation and typesetting produced by AI.
\clearpage
